\section{The complete mean-zero scalar coefficient}
\label{sgc-section}

The exact factorization of Notebook 36 reads
\[
 \det\widehat G=\det T_s\det M_s,\qquad
 T_s=D_A^TD_A+K^{-2}J_AJ_A^T.
\]
Here \(M_s=K^{-2}J_A^TT_s^{-1}J_A\) has one fixed color
fiber per final block. Its actual Maxwell coefficient is uniformly
bounded. The normalization is
\(\beta(F)=2[t^2]F\) on the specified one-color affine probe;
normalization at the flat field contributes zero.

\begin{theorem}[Mean-zero scalar logarithm]
\label{sgc-main}
Under the actual structural and affine-identification hypotheses of
Notebooks 35--36 and Theorem~\ref{spb-main},
\begin{equation}
 \left|\beta(\mathcal N_0\log\det\widehat G)
                 -2\alpha_g\log K\right|\le C,
 \qquad \alpha_g=\frac{C_H}{48\pi^2}.
 \label{sgc-coefficient}
\end{equation}
The bound is uniform in preceding depth and in every applicable
joint infinite structural system.
\end{theorem}

\begin{proof}
The scalar part of the infrared-tail theorem in Notebook 36 bounds
the integral over \(s>K^2\). Theorem~\ref{spb-main} removes the scalar
soft penalty on \(0<s\le K^2\) with a bounded integrated error.
The coarse factor \(M_s\) contributes a further bounded error.
The exact unpinned scalar heat calculation of Notebook 34 therefore
gives
\begin{equation}
 \beta(\mathcal N_0\log\det\widehat G)
       =\int_0^{K^2}h_g(s)\,\frac{ds}{s}+O(1),
 \qquad
 h_g(s)=\frac{C_H}{3}s^2e^{-8s}I_0(2s)^2I_1(2s)^2.
 \label{sgc-cutoff}
\end{equation}
The definition of \(h_g\) includes the factor \(-2\) multiplying
the scalar heat coefficient. The minus sign comes from the heat
representation of the logarithm, and the factor two from \(\beta\).
No additional real/complex color multiplicity is introduced.

The previously established elementary Bessel bounds imply
\[
 h_g(s)=O(s^4)\quad(0<s\le1),\qquad
 h_g(s)=\alpha_g+O(s^{-1})\quad(s\ge1).
 \]
Thus \(\int_0^1|h_g(s)|\,ds/s\) and
\(\int_1^\infty|h_g(s)-\alpha_g|\,ds/s\) are finite.
Splitting \eqref{sgc-cutoff} at one proves
\(\alpha_g\log(K^2)+O(1)\), including \(K=1\).
\end{proof}

This is the scalar determinant in the exact covariant mean-zero
decomposition. The earlier globally point-pinned scalar determinant
has a different logarithmic coefficient; its compensating vector
factor is different as well. Only each complete, exactly normalized
determinant combination is invariant under those changes of gauge
decomposition. This theorem alone does not compute the vector
coefficient or prove positive total drift.
